% Guide d'utilisation de l'administration — Fils Naturel Tech-Com
% Compilation : pdflatex guide-admin-fntc.tex (deux passes pour la table des matières)

\documentclass[11pt,a4paper]{article}

\usepackage{lmodern} % polices vectorielles : le PDF reste net à l'écran et à l'impression
\usepackage[T1]{fontenc}
\usepackage[utf8]{inputenc}
\usepackage[french]{babel}
\usepackage[a4paper,margin=2.4cm,headheight=15pt]{geometry}
\usepackage{xcolor}
\usepackage{titlesec}
\usepackage{enumitem}
\usepackage{fancyhdr}
\usepackage{tabularx}
\usepackage{booktabs}
\usepackage{tcolorbox}
\usepackage{graphicx}
\usepackage[hidelinks]{hyperref}

% Couleurs reprises du site : os, rouge signature, ambre.
\definecolor{sig}{HTML}{D8161F}
\definecolor{amb}{HTML}{E8B33A}
\definecolor{encre}{HTML}{1A0D0C}
\definecolor{encre2}{HTML}{6B5D57}
\definecolor{os}{HTML}{F4F2EC}
\definecolor{os2}{HTML}{EAE7DE}

\hypersetup{
  pdftitle={Guide de l'administration — Fils Naturel Tech-Com},
  pdfauthor={Fils Naturel Tech-Com},
  pdfsubject={Mode d'emploi de la partie administration},
  colorlinks=false,
}

\pagestyle{fancy}
\fancyhf{}
\fancyhead[L]{\small\color{encre2}Fils Naturel Tech-Com}
\fancyhead[R]{\small\color{encre2}Guide de l'administration}
\fancyfoot[C]{\small\color{encre2}\thepage}
\renewcommand{\headrulewidth}{0.4pt}
\renewcommand{\headrule}{\hbox to\headwidth{\color{os2}\leaders\hrule height \headrulewidth\hfill}}

\titleformat{\section}{\Large\bfseries\color{encre}}{\thesection.}{0.6em}{}
\titleformat{\subsection}{\large\bfseries\color{encre}}{\thesubsection}{0.6em}{}
\titlespacing*{\section}{0pt}{1.6em}{0.7em}

\setlist[itemize]{leftmargin=1.2em,itemsep=0.25em,topsep=0.35em}
\setlist[enumerate]{leftmargin=1.4em,itemsep=0.3em,topsep=0.35em}

% Encadrés : conseil, attention, à retenir.
\newtcolorbox{conseil}[1][]{colback=os,colframe=amb,boxrule=0pt,leftrule=3pt,
  arc=2pt,left=8pt,right=8pt,top=6pt,bottom=6pt,#1}
\newtcolorbox{attention}[1][]{colback=os,colframe=sig,boxrule=0pt,leftrule=3pt,
  arc=2pt,left=8pt,right=8pt,top=6pt,bottom=6pt,#1}

\newcommand{\bouton}[1]{\textbf{\textcolor{sig}{#1}}}
\newcommand{\menu}[1]{\textsf{\textbf{#1}}}
\newcommand{\champ}[1]{\textit{#1}}

\begin{document}

% ------------------------------------------------------------------ couverture
\begin{titlepage}
\centering
\vspace*{3cm}
{\color{sig}\rule{4cm}{2pt}}\par
\vspace{1.2cm}
{\Huge\bfseries Guide de l'administration\par}
\vspace{0.8cm}
{\LARGE Fils Naturel Tech-Com\par}
\vspace{2.5cm}
{\large Comment gérer la boutique au quotidien :\\
articles, photos, prix d'achat et de vente, stock,\\
commandes, achats et bilan des bénéfices.\par}
\vfill
{\color{encre2}\large Version 1 — \today\par}
\vspace{0.5cm}
{\color{sig}\rule{4cm}{2pt}}\par
\end{titlepage}

\tableofcontents
\newpage

% ------------------------------------------------------------------ démarrage
\section{Avant de commencer}

\subsection{Se connecter}

L'administration se trouve à l'adresse du site suivie de \texttt{/admin}. Par
exemple :

\begin{center}\ttfamily https://votre-site.onrender.com/admin\end{center}

\begin{enumerate}
  \item Saisir l'identifiant et le mot de passe reçus.
  \item Le bouton en forme d'\oe il, à droite du champ, affiche le mot de passe
        pour vérifier la saisie.
  \item Cliquer sur \bouton{Se connecter}.
\end{enumerate}

\begin{attention}
Après cinq mots de passe faux, le compte est bloqué pendant quinze minutes.
Ce n'est pas une panne : c'est la protection contre quelqu'un qui essaierait
de deviner le mot de passe.
\end{attention}

\begin{conseil}
Changez le mot de passe fourni au départ dès la première connexion, dans
\menu{Administration $\rightarrow$ Comptes admin}. Un mot de passe long
(une phrase entière) vaut mieux qu'un mot compliqué et court.
\end{conseil}

\subsection{Les quatre niveaux de compte}

Chaque personne à qui vous donnez un accès reçoit un rôle. Le rôle décide de
ce qu'elle peut faire ; le site le vérifie à chaque action, pas seulement dans
le menu.

\begin{center}
\begin{tabularx}{\textwidth}{@{}lX@{}}
\toprule
\textbf{Rôle} & \textbf{Ce qu'il permet} \\
\midrule
Lecteur & Regarder, sans rien modifier. Pour montrer les chiffres à quelqu'un. \\
Vendeur & Commandes, stock, achats, clients, messages. Le rôle du comptoir. \\
Gestionnaire & Tout ce qui précède, plus les articles, les prix, les promotions et les textes du site. \\
Propriétaire & Tout, y compris les comptes, les devises, les paramètres et le journal. \\
\bottomrule
\end{tabularx}
\end{center}

\begin{conseil}
Donnez le rôle \textit{Vendeur} à un employé du comptoir. Gardez
\textit{Propriétaire} pour vous. Le dernier propriétaire ne peut être ni
supprimé ni rétrogradé : le site refuse de vous enfermer dehors.
\end{conseil}

\subsection{Se repérer dans le menu}

Le menu de gauche est rangé par métier :

\begin{itemize}
  \item \menu{Pilotage} — le tableau de bord et le bilan des bénéfices.
  \item \menu{Catalogue} — les articles, les rayons, le stock, les achats, les promotions.
  \item \menu{Ventes} — les commandes, les clients, les messages.
  \item \menu{Contenus} — les textes du site, la médiathèque, les bannières.
  \item \menu{Administration} — les paramètres, les devises, les comptes, le journal.
\end{itemize}

Sur téléphone, le menu se replie ; tout reste accessible.

% ------------------------------------------------------------------ articles
\section{Gérer les articles}

\subsection{Créer un article}

\menu{Catalogue $\rightarrow$ Produits}, puis \bouton{Nouveau produit}.

\begin{enumerate}
  \item \champ{Nom} — celui que verront les clients.
  \item \champ{Rayon} — obligatoire ; créez-le d'abord dans \menu{Catégories} s'il manque.
  \item \champ{Référence} — votre code interne, facultatif mais pratique pour la recherche.
  \item \champ{Description courte} — la phrase qui s'affiche sous le nom dans le catalogue.
  \item \champ{Description complète} — le texte de la fiche.
  \item \champ{Caractéristiques} — les lignes du tableau technique (poids, garantie, puissance\ldots).
  \item Les prix et le stock : voir ci-dessous.
  \item \champ{Statut} — \textit{Publié} pour le mettre en ligne, \textit{Brouillon} pour le préparer sans le montrer.
  \item \bouton{Enregistrer}.
\end{enumerate}

\subsection{Prix d'achat et prix de vente}

Deux champs distincts, à ne pas confondre :

\begin{center}
\begin{tabularx}{\textwidth}{@{}lX@{}}
\toprule
\champ{Prix d'achat, l'unité} & Ce que l'article vous coûte. Sert à calculer votre bénéfice. Le client ne le voit jamais. \\
\champ{Prix de vente} & Ce que le client paie. C'est le prix affiché sur le site. \\
\bottomrule
\end{tabularx}
\end{center}

Le prix d'achat se remplit tout seul dès que vous enregistrez un achat
(chapitre~\ref{sec:achats}) : inutile de le tenir à jour à la main.

\begin{attention}
Si le prix d'achat reste à zéro, le bilan comptera l'article comme du bénéfice
entier — vos chiffres seront faux, trop beaux. La page \menu{Bilan} vous
prévient quand des articles en stock sont dans ce cas.
\end{attention}

\subsection{Ajouter des photos}

Dans le formulaire de l'article, zone \champ{Photos} : cliquez, choisissez une
ou plusieurs images, jusqu'à huit. Elles sont réduites et compressées
automatiquement avant l'envoi ; une photo prise au téléphone passe sans
problème. La première image est celle qui apparaît dans le catalogue ;
faites-la glisser pour changer l'ordre.

\subsection{Mettre un article en promotion}

Sur la fiche de l'article : \champ{Prix promo}, puis \champ{Promo à partir du}
et \champ{Promo jusqu'au}. Hors de ces dates, le prix normal revient tout seul.
Le prix barré s'affiche à côté du prix promotionnel.

Pour une réduction sur toute la commande, utilisez plutôt un code promo :
\menu{Catalogue $\rightarrow$ Promotions}.

% ------------------------------------------------------------------ stock
\section{Le stock}

\menu{Catalogue $\rightarrow$ Stock et ruptures} liste les articles du plus bas
au plus haut, les ruptures en tête. Deux façons de corriger une ligne :

\begin{itemize}
  \item \champ{Nouvelle quantité} — vous tapez le stock réel compté sur l'étagère. Utilisez-la après un inventaire.
  \item \champ{Ajouter / retirer} — vous tapez l'écart : \texttt{-2} pour deux articles cassés, \texttt{+5} pour cinq retrouvés.
\end{itemize}

Chaque correction laisse une trace datée et signée dans l'historique, en bas de
la page. Rien ne se modifie en silence.

\begin{conseil}
Pour une entrée de marchandise, ne passez pas par cette page : enregistrez un
\textbf{achat}. Le stock monte pareil, mais le prix payé est gardé et votre
bénéfice reste juste.
\end{conseil}

% ------------------------------------------------------------------ achats
\section{Les achats}
\label{sec:achats}

\menu{Catalogue $\rightarrow$ Achats}. C'est ici qu'on note ce qu'on a acheté
pour revendre.

\subsection{Enregistrer un achat}

\begin{enumerate}
  \item \champ{Article} — celui qui entre en stock.
  \item \champ{Acheté} — \textit{À l'unité} ou \textit{Par lot / carton}.
  \item \champ{Nombre} — le nombre d'unités, ou le nombre de lots.
  \item \champ{Unités par lot} — apparaît en mode lot : combien de pièces contient un carton.
  \item \champ{Le prix saisi est} — \textit{Le total payé}, ou \textit{Le prix d'un lot} / \textit{d'une unité}.
  \item \champ{Prix payé} — le montant.
  \item \champ{Date}, \champ{Fournisseur}, \champ{Note} — la date compte pour le bilan ; le reste est facultatif.
  \item Vérifiez l'encadré de contrôle, puis \bouton{Enregistrer l'achat}.
\end{enumerate}

L'encadré calcule sous vos yeux, avant validation : combien d'unités entrent,
ce que revient chaque unité, le total payé et le stock après l'opération.

\begin{conseil}
\textbf{Exemple.} Dix cartons de douze câbles, à 6\,000 le carton. Choisissez
\textit{Par lot}, nombre \texttt{10}, unités par lot \texttt{12}, prix saisi
\textit{Le prix d'un lot}, prix \texttt{6000}. Le site comprend : 120 câbles
entrent, chacun revient à 500, vous avez payé 60\,000.
\end{conseil}

\subsection{Ce que l'achat déclenche}

\begin{itemize}
  \item Le stock de l'article monte du nombre d'unités.
  \item Le prix d'achat de l'article est recalculé en \textbf{moyenne pondérée}.
        Si vous aviez 10 pièces à 400 et que 10 entrent à 600, le coût passe à 500.
  \item La dépense est datée et comptée dans le bilan de la période.
  \item Un mouvement de stock est tracé, avec votre nom.
\end{itemize}

\subsection{Corriger une erreur}

Le bouton \bouton{Supprimer} sur la ligne d'un achat annule l'opération : les
unités concernées repartent du stock. Une confirmation est demandée.

\begin{attention}
La suppression retire les unités du stock mais ne rétablit pas l'ancien prix
d'achat moyen. Après une grosse erreur, vérifiez le champ \champ{Prix d'achat}
sur la fiche de l'article et corrigez-le à la main si besoin.
\end{attention}

% ------------------------------------------------------------------ commandes
\section{Les commandes}

\subsection{Voir ce qui a été commandé}

\menu{Ventes $\rightarrow$ Commandes}. Le tableau donne, pour chaque commande,
le numéro, le client, la date, \textbf{le détail des articles} (par exemple
\texttt{2 $\times$ Routeur Wi-Fi 6}), l'état du paiement, l'état du traitement
et le total. Le bouton \bouton{Ouvrir} affiche la fiche complète : prix
unitaires, quantités, sous-total, remise, livraison, adresse et historique.

Les champs de recherche du haut filtrent par numéro, nom, email, état ou
paiement.

\subsection{Faire avancer une commande}

Sur la fiche, l'encadré \menu{Où en est cette commande} donne les boutons :

\begin{center}
\begin{tabularx}{\textwidth}{@{}lX@{}}
\toprule
\bouton{Marquer préparée} & Les articles sont rassemblés, prêts à partir. \\
\bouton{Marquer expédiée} & La commande est partie. Notez le numéro de suivi juste en dessous. \\
\bouton{Le client a reçu la commande} & La commande est livrée. \\
\bouton{Marquer payée} & L'argent est encaissé. \\
\bouton{Annuler la commande} & Voir ci-dessous. \\
\bottomrule
\end{tabularx}
\end{center}

Un bouton disparaît quand il n'a plus de sens : une commande déjà livrée ne
propose plus « Le client a reçu la commande ».

\begin{attention}
\textbf{Une vente n'entre dans le bilan que lorsqu'elle est à la fois payée et
livrée.} Tant que ces deux boutons n'ont pas été cliqués, la vente n'apparaît
ni dans le chiffre d'affaires ni dans le bénéfice.
\end{attention}

\subsection{Annuler une commande}

\bouton{Annuler la commande}, en rouge. Une confirmation est demandée, puis
\textbf{les articles retournent en stock} automatiquement, avec un mouvement
tracé. Pour rouvrir une commande annulée par erreur, utilisez le sélecteur
\menu{Changer l'état à la main}, juste en dessous des boutons.

\subsection{Saisir une commande au comptoir}

\menu{Ventes $\rightarrow$ Commandes}, bouton \bouton{Saisir une commande}.
Utile pour une vente en boutique, par téléphone ou par message.

\begin{enumerate}
  \item Cherchez les articles, filtrez par rayon si la liste est longue.
  \item Ajoutez-les et \textbf{tapez la quantité} au clavier — pas besoin de cliquer trente fois.
  \item Choisissez le client s'il est déjà inscrit : son nom, son email, son téléphone et son adresse se remplissent seuls.
  \item Ajustez les frais de livraison, le mode et l'état du paiement.
  \item \bouton{Enregistrer la commande}. Le stock baisse aussitôt.
\end{enumerate}

% ------------------------------------------------------------------ bilan
\section{Le bilan et les bénéfices}

\menu{Pilotage $\rightarrow$ Bilan et bénéfices}. Quatre onglets en haut :
\bouton{Aujourd'hui}, \bouton{Cette semaine}, \bouton{Ce mois},
\bouton{Cette année}. La semaine commence le lundi.

\subsection{Lire les chiffres}

\begin{center}
\begin{tabularx}{\textwidth}{@{}lX@{}}
\toprule
\textbf{Ventes} & Ce que les clients ont payé sur la période. \\
\textbf{Bénéfice} & Les ventes moins ce que la marchandise vous a coûté. C'est ce qui vous reste. \\
\textbf{Taux de marge} & Le bénéfice en pourcentage des ventes. 30~\% veut dire : sur 100 encaissés, 30 vous restent. \\
\textbf{Achats de la période} & Ce que vous avez dépensé en marchandise pendant la période. \\
\textbf{Unités achetées} & Le nombre de pièces entrées en stock. \\
\textbf{Reste en stock} & Le nombre de pièces disponibles, avec les ruptures et les seuils d'alerte. \\
\textbf{Stock au prix d'achat} & L'argent qui dort sur vos étagères. \\
\textbf{Stock au prix de vente} & Ce que ce stock rapportera une fois vendu, et le bénéfice à venir. \\
\bottomrule
\end{tabularx}
\end{center}

Sous chaque chiffre, la comparaison avec la période précédente : hier pour
aujourd'hui, la semaine dernière pour cette semaine, et ainsi de suite. En
vert, ça va dans le bon sens ; en rouge, ça recule.

Le tableau du bas classe vos dix meilleurs articles de la période : unités
vendues, ventes et bénéfice.

\begin{conseil}
Un bénéfice qui paraît trop beau vient presque toujours d'articles sans prix
d'achat. Le message d'alerte en haut de la page vous dit combien d'articles
sont concernés.
\end{conseil}

% ------------------------------------------------------------------ le reste
\section{Le reste de l'administration}

\subsection{Clients}
\menu{Ventes $\rightarrow$ Clients} : la fiche de chaque client, ses commandes,
son total dépensé et son badge de fidélité. La fiche est modifiable.

\subsection{Messages}
\menu{Ventes $\rightarrow$ Messages} : les conversations ouvertes depuis la
bulle de chat du site. La pastille du menu compte les messages non lus.
Répondez depuis la conversation ; le visiteur voit la réponse sans rechargement.

\subsection{Textes du site, bannières, médiathèque}
\menu{Contenus} : tous les textes de l'accueil, les bandeaux d'annonce et les
images réutilisables se modifient sans toucher au code.

\subsection{Devises et taux}
\menu{Administration $\rightarrow$ Devises et taux} : la devise de base (celle
dans laquelle vous saisissez les prix) et le taux de chaque autre devise. Le
visiteur choisit la sienne en haut du site ; les prix sont convertis à
l'affichage. Mettez les taux à jour quand ils bougent.

\subsection{Paramètres}
\menu{Administration $\rightarrow$ Paramètres} : nom du site, slogan, contacts,
frais de livraison, seuil de livraison offerte.

\subsection{Comptes admin}
\menu{Administration $\rightarrow$ Comptes admin} : créer un accès, changer un
rôle, réinitialiser un mot de passe.

\subsection{Journal d'audit}
\menu{Administration $\rightarrow$ Journal d'audit} : qui a fait quoi, et quand.
Il ne peut pas être effacé depuis l'interface. C'est votre preuve en cas de
doute sur une modification.

% ------------------------------------------------------------------ habitudes
\section{Les bonnes habitudes}

\begin{enumerate}
  \item \textbf{Chaque matin} — ouvrir \menu{Commandes}, traiter les nouvelles, marquer préparées celles qui partent.
  \item \textbf{À chaque arrivage} — enregistrer un \textbf{achat}, jamais une correction de stock.
  \item \textbf{À chaque vente livrée et payée} — cliquer les deux boutons, sinon le bilan reste faux.
  \item \textbf{Chaque semaine} — regarder \menu{Bilan} en vue hebdomadaire et la page \menu{Stock} pour les ruptures.
  \item \textbf{Chaque mois} — inventaire physique, correction par \textit{Nouvelle quantité}, puis bilan mensuel.
\end{enumerate}

% ------------------------------------------------------------------ dépannage
\section{En cas de problème}

\begin{center}
\begin{tabularx}{\textwidth}{@{}lX@{}}
\toprule
\textbf{Symptôme} & \textbf{Que faire} \\
\midrule
Le site met vingt secondes à s'ouvrir & Sur l'offre gratuite, le site s'endort après quinze minutes sans visite. La première visite le réveille. Passer à l'offre payante supprime l'attente. \\
« Compte bloqué » à la connexion & Cinq mots de passe faux. Attendre quinze minutes, ou faire réinitialiser par un propriétaire. \\
Une photo refuse de s'enregistrer & Essayer avec moins d'images à la fois, ou une image moins lourde. \\
Le stock ne correspond pas & Ouvrir \menu{Stock}, lire l'historique en bas : chaque mouvement est daté et signé. \\
Le bénéfice paraît faux & Des articles n'ont pas de prix d'achat. Voir l'alerte en haut du \menu{Bilan}. \\
Une commande a disparu de la liste & Un filtre est actif. Vider la recherche et remettre « Tous les états ». \\
\bottomrule
\end{tabularx}
\end{center}

\vspace{1.5cm}
\begin{center}
{\color{sig}\rule{3cm}{1.5pt}}\par
\vspace{0.4cm}
{\color{encre2}\small Fils Naturel Tech-Com — guide de l'administration}
\end{center}

\end{document}
